% GENERATED FILE. Edit canonical lessons, then run npm run generate:books.
% canonical-source-hash: fnv1a64:2cd8920151f02bf5
% canonical-lessons: RU-C141-vstrechat, RU-C141-znakomitsya, RU-C141-priglashat, RU-C141-blagodarit, RU-C141-izvinyatsya

\chapter{To meet, To get acquainted, To invite, To thank, To apologise}
\label{ch:ru-to-meet-to-get-acquainted-to-invite-to-thank-to-apologise}
\hlchaptermodality{141}

\begin{chapteropening}
\textbf{By the end of this chapter:} \emph{I can say to meet, to get acquainted, to invite, to thank and to apologise.}

Complete the last lesson of chapter 141: I can say to meet, to get acquainted, to invite, to thank and to apologise.
\end{chapteropening}

\section[\texorpdfstring{\ru{встречать} (vstrechát)}{встречать (vstrechát)}]{\ru{встречать} (vstrechát) — to meet}

\label{lesson:RU-C141-vstrechat}



\begin{quote}
Before the new one: say the Russian for to do the laundry, then the Russian for to iron; to stroke.
\end{quote}

\subsection*{You'll want to know: \ru{встречать}}
\textbf{\ru{встречать}} (\emph{vstrechát}) — \textquotedblleft{}to meet\textquotedblright{}.

\begin{grammarlens}[title={What you've built}]
One more word for this chapter.
\end{grammarlens}

\subsection*{Your turn}
\begin{itemize}
\raggedright
  \item \emph{Say it:} \emph{vstrechát}
  \item \emph{Say it:} \emph{vstrechát}, once more
  \item \emph{Say it:} say \emph{vstrechát}
  \item \emph{Recall:} say the Russian for to do the laundry, then the Russian for to iron; to stroke, then say \emph{vstrechát} again
\end{itemize}

\begin{tcolorbox}[breakable,colback=teal!4,colframe=teal!35!black,title={Before you move on}]
What does \ru{встречать} mean? (\textquotedblleft{}To meet\textquotedblright{}.) Say it once more.
\end{tcolorbox}
\section[\texorpdfstring{\ru{знакомиться} (znakómitsya)}{знакомиться (znakómitsya)}]{\ru{знакомиться} (znakómitsya) — to get acquainted}

\label{lesson:RU-C141-znakomitsya}



\begin{quote}
Before the new one: say the Russian for to iron; to stroke, then the Russian for to meet.
\end{quote}

\subsection*{You'll want to know: \ru{знакомиться}}
\textbf{\ru{знакомиться}} (\emph{znakómitsya}) — \textquotedblleft{}to get acquainted\textquotedblright{}.

\begin{grammarlens}[title={What you've built}]
One more word for this chapter.
\end{grammarlens}

\subsection*{Your turn}
\begin{itemize}
\raggedright
  \item \emph{Say it:} \emph{znakómitsya}
  \item \emph{Say it:} \emph{znakómitsya}, once more
  \item \emph{Say it:} say \emph{znakómitsya}
  \item \emph{Recall:} say the Russian for to iron; to stroke, then the Russian for to meet, then say \emph{znakómitsya} again
\end{itemize}

\begin{tcolorbox}[breakable,colback=teal!4,colframe=teal!35!black,title={Before you move on}]
What does \ru{знакомиться} mean? (\textquotedblleft{}To get acquainted\textquotedblright{}.) Say it once more.
\end{tcolorbox}
\section[\texorpdfstring{\ru{приглашать} (priglashát)}{приглашать (priglashát)}]{\ru{приглашать} (priglashát) — to invite}

\label{lesson:RU-C141-priglashat}



\begin{quote}
Before the new one: say the Russian for to meet, then the Russian for to get acquainted.
\end{quote}

\subsection*{You'll want to know: \ru{приглашать}}
\textbf{\ru{приглашать}} (\emph{priglashát}) — \textquotedblleft{}to invite\textquotedblright{}.

\begin{grammarlens}[title={What you've built}]
One more word for this chapter.
\end{grammarlens}

\subsection*{Your turn}
\begin{itemize}
\raggedright
  \item \emph{Say it:} \emph{priglashát}
  \item \emph{Say it:} \emph{priglashát}, once more
  \item \emph{Say it:} say \emph{priglashát}
  \item \emph{Recall:} say the Russian for to meet, then the Russian for to get acquainted, then say \emph{priglashát} again
\end{itemize}

\begin{tcolorbox}[breakable,colback=teal!4,colframe=teal!35!black,title={Before you move on}]
What does \ru{приглашать} mean? (\textquotedblleft{}To invite\textquotedblright{}.) Say it once more.
\end{tcolorbox}
\section[\texorpdfstring{\ru{благодарить} (blagodarít)}{благодарить (blagodarít)}]{\ru{благодарить} (blagodarít) — to thank}

\label{lesson:RU-C141-blagodarit}



\begin{quote}
Before the new one: say the Russian for to get acquainted, then the Russian for to invite.
\end{quote}

\subsection*{You'll want to know: \ru{благодарить}}
\textbf{\ru{благодарить}} (\emph{blagodarít}) — \textquotedblleft{}to thank\textquotedblright{}.

\begin{grammarlens}[title={What you've built}]
One more word for this chapter.
\end{grammarlens}

\subsection*{Your turn}
\begin{itemize}
\raggedright
  \item \emph{Say it:} \emph{blagodarít}
  \item \emph{Say it:} \emph{blagodarít}, once more
  \item \emph{Say it:} say \emph{blagodarít}
  \item \emph{Recall:} say the Russian for to get acquainted, then the Russian for to invite, then say \emph{blagodarít} again
\end{itemize}

\begin{tcolorbox}[breakable,colback=teal!4,colframe=teal!35!black,title={Before you move on}]
What does \ru{благодарить} mean? (\textquotedblleft{}To thank\textquotedblright{}.) Say it once more.
\end{tcolorbox}
\section[\texorpdfstring{\ru{извиняться} (izvinyátsya)}{извиняться (izvinyátsya)}]{\ru{извиняться} (izvinyátsya) — to apologise}

\label{lesson:RU-C141-izvinyatsya}



\begin{quote}
Before the new one: say the Russian for to invite, then the Russian for to thank.
\end{quote}

\subsection*{You'll want to know: \ru{извиняться}}
\textbf{\ru{извиняться}} (\emph{izvinyátsya}) — \textquotedblleft{}to apologise\textquotedblright{}.

\begin{grammarlens}[title={What you've built}]
One more word for this chapter.
\end{grammarlens}

\subsection*{Your turn}
\begin{itemize}
\raggedright
  \item \emph{Say it:} \emph{izvinyátsya}
  \item \emph{Say it:} \emph{izvinyátsya}, once more
  \item \emph{Say it:} say \emph{izvinyátsya}
  \item \emph{Recall:} say the Russian for to invite, then the Russian for to thank, then say \emph{izvinyátsya} again
\end{itemize}

\begin{tcolorbox}[breakable,colback=teal!4,colframe=teal!35!black,title={Before you move on}]
What does \ru{извиняться} mean? (\textquotedblleft{}To apologise\textquotedblright{}.) Say it once more.
\end{tcolorbox}
